\subsection{The optimal size of one die}\label{individual-section}

The arithmetic lower bound forces almost all dice to be exponentially
large, but it leaves room for a prescribed die to be much smaller.  The
following theorem determines the optimal order of its size.

\begin{theorem}[The size of one die]\label{thm:individual-size}
Every die in a permutation-fair family of \(n\geq2\) dice has at least
\((n-1)^2/200\) faces.  Conversely, there is an absolute constant \(C>0\)
such that, for every \(n\geq2\) and every integer \(K\geq Cn^2\), there
is a permutation-fair family in which a specified die has exactly \(K\)
faces.  The remaining \(n-1\) dice can all have the same finite number of
faces.
\end{theorem}

Thus, if \(g(n)\) is the smallest possible number of faces of one
specified die in a permutation-fair family of \(n\) dice, then
\[
  g(n)=\Theta(n^2).
\]
The construction below gives no quantitative bound for the other dice,
and hence no improvement in the minimum total size.

\subsubsection{A universal quadratic lower bound}

Fix a die \(D_0\) with sorted faces \(x_1<\cdots<x_K\), and put
\(m=n-1\).  For each other die \(D_i\), define
\[
  p_i(q)=\Pr(D_i<x_q),\qquad 1\leq i\leq m,\quad 1\leq q\leq K.
\]
Each \(p_i(q)\) is nondecreasing in \(q\).  Independence, followed by
permutation fairness restricted to \(D_0\) and any subset of the other
dice, gives
\begin{equation}\label{individual-comparison-moments}
  \frac1K\sum_{q=1}^K\prod_{i\in S}p_i(q)=\frac1{|S|+1}
  \qquad(S\subseteq[m]).
\end{equation}
Indeed, the left side is the probability that \(D_0\) beats every die
indexed by \(S\).  We will use these identities with general positive
weights \(w_q\) summing to one in place of \(1/K\).  By linearity,
their weighted version implies
\begin{equation}\label{individual-diagonal-identity}
  \sum_{q=1}^K w_q F(p_1(q),\ldots,p_m(q))
  =\int_0^1 F(t,\ldots,t)\,dt
\end{equation}
for every multi-affine polynomial \(F\), meaning one whose degree in
each variable is at most one.  We write \(\mathbb E_w\) for the weighted
sum on the left.

\begin{lemma}[Paired thresholds]\label{individual-paired-thresholds}
Suppose the vectors \(p(q)\in[0,1]^m\) are coordinatewise
nondecreasing in \(q\).  For distinct indices \(j,k\) and any
\(t\in(0,1)\), there are thresholds \(\alpha,\beta\) such that
\[
  \{\alpha,\beta\}=\{t,b\},\qquad 0\leq b\leq t,
  \qquad
  (p_j(q)-\alpha)(p_k(q)-\beta)\geq0
  \quad\text{for every }q.
\]
\end{lemma}

\begin{proof}
Adjoin the all-zero and all-one vectors and interpolate successive
vectors linearly to form a continuous coordinatewise nondecreasing path.
Stop at the first point where either coordinate \(j\) or coordinate
\(k\) reaches \(t\), and let \(\alpha,\beta\) be their values there.
One equals \(t\), and the other is at most \(t\).  Every original row
before this point lies below both thresholds, and every row after it lies
above both.  Thus the two differences have the same weak sign.  The
interpolating points are used only to choose thresholds; they carry no
weight in~\eqref{individual-diagonal-identity}.
\end{proof}

We need the following elementary facts about quadrature for the uniform
measure on \([0,1]\).  Their proofs are included to specify the endpoint
estimates used below.

\begin{lemma}[Gaussian and Radau rules]\label{individual-quadrature-facts}
The \(D\)-node Gaussian rule has positive weights, is exact through
degree \(2D-1\), and has largest node \(\xi_D\) satisfying
\begin{equation}\label{individual-gaussian-edge}
  1-\xi_D\leq\frac{\pi^2}{4(D+1)^2}.
\end{equation}
The right Radau rule with interior nodes \(t_1<\cdots<t_e\) and endpoint
\(1\) has positive weights and is exact through degree \(2e\).  Its
endpoint weight and interior nodes satisfy
\begin{equation}\label{individual-radau-identities}
  \omega_{\mathrm{end}}=\frac1{(e+1)^2},
  \qquad
  \sum_{\ell=1}^e\frac1{1-t_\ell}=\frac{e(e+2)}2.
\end{equation}
\end{lemma}

\begin{proof}
Let \(P_r\) denote the Legendre polynomial normalized by \(P_r(1)=1\).
The Gaussian nodes are the shifted roots of \(P_D\).  Polynomial
division by the node polynomial and orthogonality give exactness through
degree \(2D-1\); applying the rule to squared cardinal polynomials gives
positive weights.  The unshifted nodes are the eigenvalues of the
symmetric \(D\)-by-\(D\) Jacobi matrix with zero diagonal and
off-diagonal entries
\[
  \frac{k}{\sqrt{(2k-1)(2k+1)}}\geq\frac12,
  \qquad 1\leq k<D.
\]
The matrix with all these entries replaced by \(1/2\) has largest
eigenvalue \(\cos(\pi/(D+1))\), with a positive sine eigenvector.
Using that vector in the Rayleigh quotient shows that the largest
unshifted Gaussian node is at least \(\cos(\pi/(D+1))\).  Rescaling to
\([0,1]\) and using \(1-\cos x\leq x^2/2\) proves
\eqref{individual-gaussian-edge}.

For the Radau rule, put
\[
  J_e(t)=\sum_{r=0}^e(2r+1)P_r(2t-1).
\]
Legendre orthogonality gives the reproducing identity
\(\int_0^1J_e(t)f(t)\,dt=f(1)\) for \(\deg f\leq e\).
Consequently \(J_e\) is orthogonal to all polynomials of degree at most
\(e-1\) for the weight \(1-t\), and has \(e\) distinct roots in
\((0,1)\).  Interpolation at these roots and at \(1\) gives a rule
exact through degree \(2e\): division by \((t-1)J_e(t)\) leaves a
remainder of degree at most \(e\), and the quotient has degree at most
\(e-1\).  The integral of the quotient term vanishes by orthogonality.
Applying the rule to
\((1-t)(J_e(t)/(t-t_\ell))^2\) shows that every interior weight is
positive.  At the endpoint, exactness on \(J_e^2\) gives
\[
  \omega_{\mathrm{end}}J_e(1)^2
  =\int_0^1J_e(t)^2\,dt=J_e(1)=(e+1)^2.
\]
Finally, the Legendre endpoint derivative identity gives
\[
  \frac{J_e'(1)}{J_e(1)}
  =\frac{\sum_{r=0}^e(2r+1)r(r+1)}{(e+1)^2}
  =\frac{e(e+2)}2.
\]
The logarithmic derivative of \(J_e\) at \(1\) is
\(\sum_{\ell=1}^e(1-t_\ell)^{-1}\), proving the remaining identity.
\end{proof}

\begin{lemma}[The last row]\label{individual-last-row}
Suppose \(w_q>0\), \(\sum_qw_q=1\), the vectors \(p(q)\in[0,1]^m\)
are coordinatewise nondecreasing, and
\eqref{individual-diagonal-identity} holds.  Write \(a_i=p_i(K)\).
Then
\[
  a_i\geq1-\frac{\pi^2}{(m+1)^2}\quad(1\leq i\leq m),
  \qquad w_K\leq\frac{200}{m^2}.
\]
\end{lemma}

\begin{proof}
Fix \(i\), put \(d=\lfloor(m-1)/2\rfloor\), and choose \(d\)
disjoint pairs of coordinates other than \(i\).  Let
\(\xi_1<\cdots<\xi_D\) be the nodes of the Gaussian rule with
\(D=d+1\) nodes.  For pair \(\ell\), apply
Lemma~\ref{individual-paired-thresholds} at \(\xi_\ell\), omitting
the largest node \(\xi_D\).  Let \(T\) be the product of the resulting
paired factors.  It is multi-affine and nonnegative on every row, and
\[
  P(t):=T(t,\ldots,t)
  =\prod_{\ell=1}^d(t-\xi_\ell)(t-b_\ell),
  \qquad b_\ell\leq\xi_\ell<\xi_D.
\]
Since \(a_i-p_i(q)\geq0\) and coordinate \(i\) is absent from \(T\),
\eqref{individual-diagonal-identity} and Gaussian exactness yield
\[
  0\leq\mathbb E_w[(a_i-p_i)T(p)]
  =\int_0^1(a_i-t)P(t)\,dt
  =\omega_D(a_i-\xi_D)P(\xi_D).
\]
Here \(\omega_D>0\) and \(P(\xi_D)>0\), also when \(d=0\).
Therefore \(a_i\geq\xi_D\).  Since \(D+1\geq(m+1)/2\), we obtain
\[
  1-a_i\leq\delta:=\frac{\pi^2}{(m+1)^2}
  \qquad\text{for every }i.
\]

For \(m\geq10\), put \(e=\lfloor m/10\rfloor\) and choose \(e\)
disjoint pairs of coordinates.  Apply
Lemma~\ref{individual-paired-thresholds} at the interior nodes of the
right Radau rule, one node for each pair.  Again let \(T\) be the
product of the paired factors, and write
\[
  P(t)=T(t,\ldots,t)
  =\prod_{\ell=1}^e(t-t_\ell)(t-b_\ell),\qquad b_\ell\leq t_\ell.
\]
Its degree is \(2e\leq m\).  Since \(P\) vanishes at the interior
Radau nodes, exactness gives
\begin{equation}\label{individual-radau-test}
  \mathbb E_w T(p)=\frac{P(1)}{(e+1)^2}.
\end{equation}
Put \(S=\sum_\ell(1-t_\ell)^{-1}=e(e+2)/2\).  The inequalities
\(e\leq m/10\) and \(e+2\leq3m/10\) imply
\[
  \delta S\leq\frac{3\pi^2}{200}<\frac14.
\]
In particular \(\delta<1-t_\ell\) for every \(\ell\), so all paired
thresholds lie strictly below all \(a_i\).  Each paired factor at the
last row, divided by its value at the all-one vector, is at least
\((1-\delta/(1-t_\ell))^2\).  Using
\(\prod_j(1-z_j)\geq1-\sum_jz_j\) for \(0\leq z_j\leq1\), we get
\[
  \frac{T(a)}{P(1)}
  \geq\prod_{\ell=1}^e
       \left(1-\frac{\delta}{1-t_\ell}\right)^2
  \geq1-2\delta S>\frac12.
\]
All rows contribute nonnegatively to
\eqref{individual-radau-test}.  Hence
\(w_K\leq2/(e+1)^2\leq200/m^2\).  For \(m<10\), the same bound
follows directly from \(w_K\leq1\).
\end{proof}

\begin{proof}[Proof of the lower bound in Theorem~\ref{thm:individual-size}]
Apply Lemma~\ref{individual-last-row} to the comparison probabilities
of any chosen die, with \(w_q=1/K\) and \(m=n-1\).
Equation~\eqref{individual-comparison-moments} supplies the required
identities, and the conclusion is \(K\geq(n-1)^2/200\).
\end{proof}

\subsubsection{A construction with a prescribed number of faces}

The upper bound starts with a real equal-weight quadrature with \(K\)
nodes.  We approximate the nodes by rationals and correct the resulting
order probabilities by changing the distributions of the other dice.
The corrections have disjoint supports and vanishing moments, so their
effect on every order probability is exactly linear.  Finally, we replace
the corrected continuous distributions by finite uniform dice without
changing the number \(K\) of distinguished faces.

\begin{lemma}[A regular starting quadrature]\label{individual-real-rule}
There is an absolute constant \(C_0\) such that, for every \(m\geq1\)
and every integer \(K\geq C_0(m+1)^2\), there are distinct nodes
\(0<u_1<\cdots<u_K<1\) satisfying
\[
  \frac1K\sum_{q=1}^K p(u_q)=\int_0^1p(t)\,dt
  \qquad(\deg p\leq m).
\]
We may take \(C_0\) large enough that \(K\geq m^2\).
\end{lemma}

\begin{proof}
Theorem~1.1 of Gilboa and Peled~\cite{gilboa2017}, applied to the
constant weight, gives an equal-weight quadrature of degree \(d\) with
exactly \(K\) nodes for every integer \(K\geq C_1d^2\).  In their
notation, the theorem bounds the threshold \(\overline N_w(d)\) for
all sufficiently large cardinalities, as well as the minimum
cardinality \(N_w(d)\).  For the constant weight their local-mass
quantity is of order \(d^2\); affine rescaling gives the assertion on
\([0,1]\).

Apply this result at degree \(2m+2\), with the prescribed \(K\).
The rule has at least \(m+2\) distinct support points: otherwise the
square of its support-vanishing polynomial would have degree at most
\(2m+2\), zero quadrature sum, and positive integral.  At least \(m\)
support points therefore lie in \((0,1)\).  Choose one node at each
of \(m\) distinct interior points as a pivot.  The Jacobian of the first
\(m\) moment sums with respect to these nodes has entries
\(s u_q^{s-1}/K\), \(1\leq s\leq m\), and is an invertible scaled
Vandermonde matrix.  The analytic implicit function theorem allows
arbitrary sufficiently small perturbations of all other nodes, with
compensating changes of the pivots preserving these \(m\) moments.

The pivots remain distinct and interior.  Move each free endpoint node
slightly into \((0,1)\).  Collisions between free nodes form proper
hyperplanes in the free parameters.  If a free node initially coincides
with a pivot, varying that free node alone changes the matching pivot
with derivative \(-1\), because their moment-Jacobian columns are
equal.  The difference of the two coordinates thus has nonzero derivative
at the initial point, so their collision set is a proper analytic zero
set.  Initially distinct coordinates remain distinct in a sufficiently
small neighborhood.  A finite union of these proper zero sets cannot
cover the open set of inward perturbations.  Choose a collision-free
perturbation and sort the nodes.  All \(K\) nodes are now distinct and
interior, and the moments through degree \(m\) are unchanged.
\end{proof}

\begin{lemma}[Local moment corrections]\label{individual-centered-bumps}
For distinct interior nodes \(u_q\), choose pairwise disjoint open
intervals \(I_q\subset(0,1)\) around them.  After shrinking these
intervals if necessary, there are neighborhoods \(V_q\subset I_q\)
of \(u_q\) and bounded piecewise-constant functions \(h_{q,t}\),
\(t\in V_q\), supported in \(I_q\), such that
\begin{equation}\label{individual-bump-moments}
  \int_0^t h_{q,t}(x)g(x)\,dx=g(t),\qquad
  \int_t^1 h_{q,t}(x)g(x)\,dx=-g(t)
  \qquad(\deg g\leq m-1).
\end{equation}
Their breakpoints are rational, and their values are rational when
\(t\) is rational.  Their coefficient vectors depend continuously on
\(t\), and the functions are uniformly bounded for \(t\) near \(u_q\).
In particular, their full moments through degree \(m-1\) vanish.
\end{lemma}

\begin{proof}
Inside each \(I_q\), choose \(m\) disjoint rational-endpoint intervals
strictly to the left of \(u_q\), and another \(m\) strictly to its
right.  Choose \(V_q\) between these two groups.  For any \(m\)
ordered disjoint intervals \(J_1<\cdots<J_m\) of positive length,
the matrix
\[
  M_{ra}=\int_{J_a}x^r\,dx,
  \qquad 0\leq r\leq m-1,\quad 1\leq a\leq m,
\]
is rational and invertible.  Indeed, multilinearity of the determinant
gives
\[
  \det M=\int_{J_1\times\cdots\times J_m}
         \prod_{a<b}(x_b-x_a)\,dx_1\cdots dx_m>0.
\]
Solve the left interval system with right-hand side
\((1,t,\ldots,t^{m-1})^T\), and the right interval system with its
negative.  Use the solutions as the values of \(h_{q,t}\) on those
intervals, and set it to zero elsewhere.  The equations give
\eqref{individual-bump-moments} on the monomial basis.  The matrices
are fixed, so their solutions have all the stated rationality,
continuity, and boundedness properties.
\end{proof}

The next lemma explains why these local corrections can repair every
full ordering simultaneously.  Write \(\mathcal P_r\) for the real
polynomials of degree at most \(r\).

\begin{lemma}[The response of order probabilities]\label{individual-order-response}
Choose \(t_q\in V_q\), pairwise disjoint sets
\(Q_1,\ldots,Q_m\subseteq[K]\), and coefficients \(\delta_{j,q}\)
such that
\[
  f_j(x)=1+\sum_{q\in Q_j}\delta_{j,q}h_{q,t_q}(x),\qquad j\in[m],
\]
are nonnegative densities.  Let \(X_0,\ldots,X_m\) be independent,
with \(X_0\) uniform on \(t_1,\ldots,t_K\) and \(X_j\) having density
\(f_j\) for \(j\in[m]\).  Define linear functionals on
\(\mathcal P_{m-1}\) by
\[
  B_j(a)=\frac1K\sum_{q\in Q_j}\delta_{j,q}a(t_q).
\]
If all \(B_j\) equal the same functional \(B\), then every specified
full ordering in which \(X_0\) has exactly \(k\) predecessors has
probability
\begin{equation}\label{individual-response-formula}
  \frac1K\sum_{q=1}^K p_k(t_q)+B(p_k'),
  \qquad
  p_k(t)=\frac{t^k(1-t)^{m-k}}{k!(m-k)!}.
\end{equation}
\end{lemma}

\begin{proof}
Fix the full ordering and condition on \(X_0=t_\ell\).  Expand the
product of the densities of the other variables.  Choosing the constant
term from every density gives probability \(p_k(t_\ell)\).

Consider a term choosing bumps for \(b\geq1\) variables, and expand
their sums over the bump indices.  These indices are distinct because
the sets \(Q_j\) are disjoint.  The chosen supports are therefore
disjoint and ordered.  If their order is inconsistent with the specified
full ordering, the term is zero.  Otherwise, integrate out the variables
whose selected density contribution was the constant \(1\).  The
resulting kernel is a product of powers of gaps between successive
chosen variables, \(t_\ell\), and the endpoints \(0,1\), divided by
factorials.  On the relevant ordered region it is a polynomial of total
degree \(m-b\).

Only the support \(I_\ell\) can contain \(t_\ell\).  If \(b\geq2\),
at least one chosen bump has its whole support on one side of
\(t_\ell\).  Because the support intervals are disjoint and ordered,
varying this coordinate within its support crosses neither another
chosen coordinate nor \(t_\ell\).  Hence the same gap-polynomial formula
holds throughout that support.  For fixed values of the other chosen variables, its
integration is over its entire support against a polynomial of degree
at most \(m-b\leq m-1\).  Its vanishing full moments make the term
zero.  The same argument eliminates a one-bump term unless its index
is \(\ell\).  Thus only one-bump terms cut by \(t_\ell\) can survive.

For such a bump from a variable before \(X_0\), let \(a\) variables
precede that variable and \(b\) lie between it and \(X_0\), so
\(a+b=k-1\).  Its kernel, with \(t=t_\ell\), is
\[
  \frac{x^a(t-x)^b(1-t)^{m-k}}{a!b!(m-k)!}.
\]
The left identity in~\eqref{individual-bump-moments} evaluates this
polynomial at \(x=t\), giving zero unless \(b=0\).  Hence only the
variable immediately preceding \(X_0\) contributes, with polynomial
\[
  A_k(t)=\frac{t^{k-1}(1-t)^{m-k}}{(k-1)!(m-k)!}
  \qquad(k\geq1).
\]
For a variable after \(X_0\), the kernel is
\[
  \frac{t^k(x-t)^a(1-x)^b}{k!a!b!},\qquad a+b=m-k-1.
\]
The right identity in~\eqref{individual-bump-moments} takes its
negative value at \(x=t\).  Only the variable immediately following
\(X_0\) contributes, with the negative of
\[
  C_k(t)=\frac{t^k(1-t)^{m-k-1}}{k!(m-k-1)!}
  \qquad(k\leq m-1).
\]
Define \(A_0=0\) and \(C_m=0\).  Both polynomials have degree at most
\(m-1\), and \(A_k-C_k=p_k'\).  Averaging over \(t_\ell\) applies
the functional \(B_j\) of the predecessor to \(A_k\) and that of the
successor to \(C_k\), omitting a term when the neighbor is absent.
Since all these functionals equal \(B\), the total correction is
\(B(p_k')\), as claimed.
\end{proof}

\begin{proof}[Proof of the construction in Theorem~\ref{thm:individual-size}]
Put \(m=n-1\) and fix the prescribed integer \(K\geq Cn^2\), where
\(C\) is large enough for Lemma~\ref{individual-real-rule}.  Choose
its nodes \(u_q\) and the corresponding neighborhoods and corrections
from Lemma~\ref{individual-centered-bumps}.  As \(K\geq m^2\), we can
choose disjoint sets \(Q_1,\ldots,Q_m\), each of size \(m\).

Approximate each \(u_q\) by a rational \(t_q\in V_q\), retaining their
strict order.  For \(1\leq s\leq m\), put
\[
  \mu_s(t)=\frac1K\sum_{q=1}^K t_q^s,
  \qquad a_s(t)=\frac{1/(s+1)-\mu_s(t)}s.
\]
For each \(j\in[m]\), solve the rational Vandermonde system
\begin{equation}\label{individual-correction-system}
  \frac1K\sum_{q\in Q_j}\delta_{j,q}t_q^{s-1}=a_s(t)
  \qquad(1\leq s\leq m).
\end{equation}
The distinct nodes make its matrix invertible.  Its solution depends
continuously on all the nodes near \(u\) and tends to zero as \(t\to u\),
since the starting quadrature gives \(a_s(u)=0\).  The bumps are
uniformly bounded near these starting nodes.  A sufficiently close
rational approximation therefore ensures that
\[
  f_j(x)=1+\sum_{q\in Q_j}\delta_{j,q}h_{q,t_q}(x)\geq\frac12
  \qquad\text{for every }x\in[0,1].
\]
The full zeroth moments of the bumps vanish, so each \(f_j\) integrates
to one.  These are positive piecewise-constant densities with rational
breakpoints and rational values.

Equation~\eqref{individual-correction-system} makes all functionals
\(B_j\) in Lemma~\ref{individual-order-response} equal to one
functional \(B\).  Checking on monomials and constants shows that,
for every \(p\in\mathcal P_m\),
\[
  B(p')=\int_0^1p(t)\,dt-\frac1K\sum_{q=1}^Kp(t_q).
\]
The lemma therefore gives every full ordering probability
\[
  \int_0^1\frac{t^k(1-t)^{m-k}}{k!(m-k)!}\,dt
  =\frac1{(m+1)!}.
\]
We have obtained a permutation-fair continuous model with one uniform
\(K\)-point variable.

It remains to realize this model by finite uniform dice.  Partition
\([0,1]\) at all breakpoints of the \(f_j\) and at every \(t_q\).
Ignore the boundary points, which have probability zero for the
continuous variables.  On each remaining interval \(J_g\), all the
densities are positive rational constants.  Thus the masses
\[
  \alpha_{j,g}=\int_{J_g}f_j(x)\,dx
\]
are positive rationals with \(\sum_g\alpha_{j,g}=1\) for every \(j\).
Choose a common denominator \(D\) for these masses.  By
Theorem~\ref{thm:first-size}, take an equal-sided permutation-fair
family on the \(m\) other dice, with \(F\) faces each.  For each
interval \(J_g\), make a copy of its fair word and replace each
occurrence of letter \(j\) by \(D\alpha_{j,g}\) consecutive copies.
Claim~\ref{clm:restriction-blowup}(ii) shows that the resulting block
is permutation-fair; it has \(FD\alpha_{j,g}\) faces of die \(j\).
Concatenate these blocks in interval order and insert one distinguished
face at each boundary \(t_q\).  Every other die now has exactly \(FD\)
faces, whereas the distinguished die still has exactly \(K\).

For each other die, the probabilities of selecting the different
intervals are exactly \(\alpha_{j,g}\), as in the continuous model.
Conditional on these interval choices, variables in different intervals
have a determined relative order.  Within a common interval, the
continuous variables are independent uniforms on that interval, so their
relative order is uniform.  The corresponding discrete variables also
have a uniform relative order, by fairness of the block and its
restrictions, using Claim~\ref{clm:restriction-blowup}(i).  Different
interval groups involve disjoint sets of independent dice, so their
internal orders combine independently.  Consequently the full order
distribution, including the distinguished faces at the interval
boundaries, is unchanged.  Assigning consecutive distinct integer labels
to the concatenated word produces the required finite dice.
\end{proof}
